Federated Learning (FL) is a compelling solution for constructing a shared (global) model from several local data sources that inherently cannot be exchanged or aggregated~\citep{fed-learn, fed-learn-adv, fed-learn-survey, fed-learn-survey2, fed-learn-survey3, learnings-fl}.
Learning without sharing data becomes crucial when data privacy or security is paramount, like healthcare applications~\citep{fed-learn-app, fed-learn-health, fl-med-gan}. 
FL operates through an iterative procedure involving rounds of model improvement.
These rounds start with distributing the current global model to local entities (users) and selecting participants to contribute to the model update.
The chosen users participate by training their local copies of the model using their respective data and returning the updated model.
Next, aggregation of the returned models produces a new global model.
The described process represents the typical procedure for FL solutions, which learns only one shared model for all users.
The most widely embraced FL method is FedAvg~\citep{fedavg}, which computes the new global model as the parameter-wise average of returned local models.

When users have similar data, a single model can perform well.
However, when local data are very diverse, ensuring good performance across users with a single shared global model can be challenging \citep{single-model-problems, single-model-problems2}. 
Recent works \citep{pfl-bench, pfl} argue that alternatively, the focus should be on Personalized Federated Learning (PFL) methods that output models tailored to the local data.
However, learning additional models can require extra memory and computation, and managing multiple models in PFL may lead to heightened communication requirements between users, which can be a bottleneck in real-world applications.
We desire a method that can perform well with heterogeneous user data and has a computational cost similar to FL techniques.

\edit{
Next, meta-learning has been successfully applied to PFL to address the challenge of diverse user data without requiring additional models during training.
Instead of learning multiple models during the federated training process, meta-learning methods focus on producing a shared model that can quickly adapt to local datasets through fine-tuning.
This approach significantly reduces the computational and memory overhead associated with PFL methods.
Building on the foundational work by~\citet{reptile}, which demonstrated that meta-learning solutions inherently emphasize varying \textbf{fixed} balances between initial model success (shared model performance on local datasets) and rapid personalization (personalized model performance after fine-tuning), we propose an enhancement tailored for FL applications.
While traditional meta-learning prioritizes adaptation across a wide distribution of tasks, the context of FL involves a more constrained task distribution: the users' local datasets. 
Due to this domain-specific constraint, focusing more on initial model success can lead to more prominent personalized model performance after fine-tuning. 
We leverage within-round learning rate decay to balance initial model performance and rapid personalization \textbf{flexibly}. 
By adjusting the decay hyperparameter, our method enables application-specific emphasis, providing a customizable approach to optimize performance for varying levels of data heterogeneity.
Crucially, this modification generalizes the widely used FedAvg algorithm~\citep{fedavg}, maintaining its computational simplicity and convergence rate~\citep{fedavg-conv}, requiring no additional gradients or models to be computed or stored.
In Section~\ref{sec: decay-method-decay}, we mathematically justify how this flexible balance leads to better performance across heterogeneous data scenarios, enabling improved outcomes for new and existing users.
}

Finally, in Section~\ref{sec: decay-exp}, we provide extensive experimental results spanning diverse data sets across vision, text, and graph applications. 
We show that within-round learning rate decay enhances FL techniques and closes remaining performance gaps with PFL methods without the additional computation and communication overhead required by such methods.
Importantly, we consistently observe a 1 to 4 percentage point improvement in average test set accuracy for new and existing users over FedAvg.

\paragraph{Roadmap.}
The organization of this work is as follows.
First, Section~\ref{sec: decay-related} briefly summarizes related work from personalized federated learning and meta-learning.
Second, Section~\ref{sec: decay-method} introduces our approach, mathematically justifies its improved flexibility, and establishes its convergence rate.
Finally, Section~\ref{sec: decay-exp} gives empirical evidence for performance improvements and computational cost compared with other benchmark methods.

